\documentclass[10pt,a4paper]{amsart}
\usepackage{amsmath,amssymb,amsthm,graphicx}
\usepackage{verbatim}
\newenvironment{reference}{\comment}{\endcomment}
%\newenvironment{reference}{}{}
\newenvironment{slogan}{\comment}{\endcomment}
\newenvironment{history}{\comment}{\endcomment}

% For commutative diagrams we use Xy-pic
\usepackage[all]{xy}

% We use 2cell for 2-commutative diagrams.
\xyoption{2cell}
% Theorem environments.
%
\theoremstyle{plain}
\newtheorem{theorem}[subsection]{Theorem}
\newtheorem{proposition}[subsection]{Proposition}
\newtheorem{lemma}[subsection]{Lemma}

\theoremstyle{definition}
\newtheorem{definition}[subsection]{Definition}
\newtheorem{example}[subsection]{Example}
\newtheorem{exercise}[subsection]{Exercise}
\newtheorem{situation}[subsection]{Situation}

\theoremstyle{remark}
\newtheorem{remark}[subsection]{Remark}
\newtheorem{remarks}[subsection]{Remarks}

\numberwithin{equation}{subsection}

% Macros
%
\def\lim{\mathop{\mathrm{lim}}\nolimits}
\def\colim{\mathop{\mathrm{colim}}\nolimits}
\def\Spec{\mathop{\mathrm{Spec}}}
\def\Hom{\mathop{\mathrm{Hom}}\nolimits}
\def\Ext{\mathop{\mathrm{Ext}}\nolimits}
\def\SheafHom{\mathop{\mathcal{H}\!\mathit{om}}\nolimits}
\def\SheafExt{\mathop{\mathcal{E}\!\mathit{xt}}\nolimits}
\def\Sch{\mathit{Sch}}
\def\Mor{\mathop{\mathrm{Mor}}\nolimits}
\def\Ob{\mathop{\mathrm{Ob}}\nolimits}
\def\Sh{\mathop{\mathit{Sh}}\nolimits}
\def\NL{\mathop{N\!L}\nolimits}
\def\CH{\mathop{\mathrm{CH}}\nolimits}
\def\proetale{{pro\text{-}\acute{e}tale}}
\def\etale{{\acute{e}tale}}
\def\QCoh{\mathit{QCoh}}
\def\Ker{\mathop{\mathrm{Ker}}}
\def\Im{\mathop{\mathrm{Im}}}
\def\Coker{\mathop{\mathrm{Coker}}}
\def\Coim{\mathop{\mathrm{Coim}}}

% Boxtimes
%
\DeclareMathSymbol{\boxtimes}{\mathbin}{AMSa}{"02}

%
% Macros for moduli stacks/spaces
%
\def\QCohstack{\mathcal{QC}\!\mathit{oh}}
\def\Cohstack{\mathcal{C}\!\mathit{oh}}
\def\Spacesstack{\mathcal{S}\!\mathit{paces}}
\def\Quotfunctor{\mathrm{Quot}}
\def\Hilbfunctor{\mathrm{Hilb}}
\def\Curvesstack{\mathcal{C}\!\mathit{urves}}
\def\Polarizedstack{\mathcal{P}\!\mathit{olarized}}
\def\Complexesstack{\mathcal{C}\!\mathit{omplexes}}
% \Pic is the operator that assigns to X its picard group, usage \Pic(X)
% \Picardstack_{X/B} denotes the Picard stack of X over B
% \Picardfunctor_{X/B} denotes the Picard functor of X over B
\def\Pic{\mathop{\mathrm{Pic}}\nolimits}
\def\Picardstack{\mathcal{P}\!\mathit{ic}}
\def\Picardfunctor{\mathrm{Pic}}
\def\Deformationcategory{\mathcal{D}\!\mathit{ef}}
\title{K-trace and K-determinant of an L-linear map factor through the L-trace and L-norm}
\author{The Stacks Project Authors}
\date{stacks-project:fields.tex:local-snapshot-20261012; GNU FDL 1.2 or later}
\begin{document}
\maketitle
\noindent\textit{Editorial note.} All delivered mathematical text is the author's own, copied from the pinned Stacks Project chapter file fields.tex; no statement or proof step was rewritten, reordered or removed.\par\medskip
\noindent\textit{Editorial note.} The counted proof is the author's complete argument: The case E\_1(a) uses the block decomposition of a : L -> L over a K-basis, whose K-determinant is by definition Norm\_\{L/K\}(a) (context span).\par\medskip
\noindent\textit{Editorial note.} The label lemma-trace-and-norm-linear names the counted statement, and the proof environment immediately following it in the source is delivered verbatim as its proof.\par\medskip
\noindent\textit{Editorial note.} The counted statement is carried by the author's own lemma environment in the chapter source.\par\medskip
\noindent\textit{Allowed standard prerequisites (editorial).}\begin{itemize}
\item Context span 2596-2610: the definitions of field trace and norm as K-trace and K-determinant of the multiplication map, with K-linearity of the trace and multiplicativity of the norm
\item Linear algebra: trace is additive and det multiplicative in the matrix; GL\_n is generated by elementary matrices of the two displayed types
\end{itemize}\medskip
\begin{definition}
\label{definition-trace-norm}
Let $L/K$ be a finite extension of fields. For $\alpha \in L$ we define
the {\it trace}
$\text{Trace}_{L/K}(\alpha) = \text{Trace}_K(\alpha : L \to L)$
and the {\it norm}
$\text{Norm}_{L/K}(\alpha) = \det_K(\alpha : L \to L)$.
\end{definition}

\noindent
It is clear from the definition that
$\text{Trace}_{L/K}$ is $K$-linear and satisfies
$\text{Trace}_{L/K}(\alpha) = [L : K]\alpha$ for $\alpha \in K$.
Similarly $\text{Norm}_{L/K}$ is multiplicative and
$\text{Norm}_{L/K}(\alpha) = \alpha^{[L : K]}$ for $\alpha \in K$.

\begin{lemma}
\label{lemma-trace-and-norm-linear}
Let $L/K$ be a finite extension of fields. Let $V$ be a finite dimensional
vector space over $L$. Let $\varphi : V \to V$ be an $L$-linear map.
Then
$$
\text{Trace}_K(\varphi : V \to V) =
\text{Trace}_{L/K}(\text{Trace}_L(\varphi : V \to V))
$$
and
$$
\det\nolimits_K(\varphi : V \to V) =
\text{Norm}_{L/K}(\det\nolimits_L(\varphi : V \to V))
$$
\end{lemma}

\begin{proof}
Choose an isomorphism $V = L^{\oplus n}$ so that $\varphi$ corresponds
to an $n \times n$ matrix. In the case of traces, both sides of the formula
are additive in $\varphi$. Hence we can assume that $\varphi$
corresponds to the matrix with exactly one nonzero entry in the $(i, j)$ spot.
In this case a direct computation shows both sides are equal.

\medskip\noindent
In the case of norms both sides are zero if $\varphi$ has a nonzero kernel.
Hence we may assume $\varphi$ corresponds to an element of
$\text{GL}_n(L)$. Both sides of the formula are multiplicative in $\varphi$.
Since every element of $\text{GL}_n(L)$ is a product of elementary
matrices we may assume that $\varphi$ either looks like
$$
E_{12}(\lambda) =
\left(
\begin{matrix}
1 & \lambda & \ldots \\
0 & 1 & \ldots \\
\ldots & \ldots & \ldots
\end{matrix}
\right)
\quad\text{or}\quad
E_1(a) =
\left(
\begin{matrix}
a & 0 & \ldots \\
0 & 1 & \ldots \\
\ldots & \ldots & \ldots
\end{matrix}
\right)
$$
(because we may also permute the basis elements if we like).
In both cases the formula is easy to verify by direct computation.
\end{proof}

\end{document}
